% ============================================================================
% CHAPTER 3: Lernpfad
% ============================================================================
\renewcommand{\chaptertag}{Bildfilter}
\chapter{Lernpfad}

% --- Local color + stamp marker for this chapter and the next ---
\definecolor{cinnabar}{RGB}{177, 56, 46}
\newcommand{\stempel}[2]{%
  \par\vspace{3pt}\noindent
  \begin{tikzpicture}[baseline=-0.3em]
    \node[fill=cinnabar, rounded corners=2pt,
          inner xsep=5pt, inner ysep=2pt,
          font=\scriptsize\bfseries\color{white}] {Stempel #1};
  \end{tikzpicture}\,%
  \textit{\color{cinnabar}#2}\par\vspace{2pt}%
}

\vspace{-5pt}
{\color{maincolor}\rule{\textwidth}{0.4pt}}
\vspace{4pt}

{\color{maincolor}\textbf{Wegweiser:}}
\begin{itemize}[leftmargin=*]
  \item Jeder Schritt baut auf dem vorhergehenden auf.
  \item Jeder Schritt enth\"alt: was du tun sollst, ein kleines Code-Beispiel
    und welchen Stempel du dabei sammelst.
  \item Wenn du eine Station geschafft hast, zeig der Lehrperson dein
    Resultat, und sie setzt den Stempel direkt in dein Heft.
\end{itemize}
\vspace{-2pt}
{\color{maincolor}\rule{\textwidth}{0.4pt}}
\vspace{8pt}

% ============================================================================
\section{Schritt 0: Bild laden}

Bevor du Filter schreiben kannst, musst du ein Bild laden. Dazu
verwenden wir die Python-Bibliothek \texttt{Pillow}. Lege ein Foto in
denselben Ordner wie dein Python-Programm und l\"ade es mit folgendem
Code:

\begin{lstlisting}
from PIL import Image

img = Image.open("foto.jpg").convert("RGB")
breite, hoehe = img.size

bild = []
for y in range(hoehe):
    zeile = []
    for x in range(breite):
        zeile.append(img.getpixel((x, y)))
    bild.append(zeile)

print("Hoehe:", hoehe, " Breite:", breite)
\end{lstlisting}

Die ersten Zeilen \"offnen das Bild und holen seine Gr\"osse. Danach
bauen die zwei verschachtelten Schleifen das Bild als verschachtelte
Liste von Pixel-Tupeln $(R, G, B)$ auf, mit der du im Rest des Projekts
arbeitest. Zum Speichern eines Bildes verwendest du am Ende:

\begin{lstlisting}
ergebnis_img = Image.new("RGB", (breite, hoehe))
for y in range(hoehe):
    for x in range(breite):
        ergebnis_img.putpixel((x, y), bild[y][x])
ergebnis_img.save("ergebnis.png")
\end{lstlisting}

\stempel{1}{Erstes Bild}

% ============================================================================
\section{Schritt 1: Erste Pixel ver\"andern}

Bevor wir richtige Filter schreiben, probieren wir aus, was passiert, wenn
wir einzelne Farbkan\"ale ver\"andern.

\subsection*{Aufw\"armer: Rot weg}

Setze den Rot-Wert aller Pixel auf 0. Was passiert?

\begin{lstlisting}[style=pseudo]
Funktion kein_rot(bild):
    Erstelle ein leeres neues_bild.
    Fuer jede Zeile des Bildes:
        Erstelle eine leere neue_zeile.
        Fuer jeden Pixel (r, g, b) der Zeile:
            Haenge (0, g, b) an neue_zeile an.
        Haenge neue_zeile an neues_bild an.
    Gib neues_bild zurueck.
\end{lstlisting}

\stempel{2}{Kein Rot}

\subsection*{Filter 1: Invertieren}

Beim Invertieren wird jeder Farbwert ``umgedreht'': aus 0 wird 255, aus 255
wird 0, aus 100 wird 155.

\begin{definitionbox}{Invertieren}
F\"ur jeden Pixel $(R, G, B)$:
\[ (R', G', B') = (255 - R,\; 255 - G,\; 255 - B) \]
\end{definitionbox}

\begin{lstlisting}[style=pseudo]
Funktion invertieren(bild):
    Erstelle ein leeres neues_bild.
    Fuer jede Zeile des Bildes:
        Erstelle eine leere neue_zeile.
        Fuer jeden Pixel (r, g, b) der Zeile:
            Haenge (255-r, 255-g, 255-b) an neue_zeile an.
        Haenge neue_zeile an neues_bild an.
    Gib neues_bild zurueck.
\end{lstlisting}

\stempel{3}{Spiegelwelt}

% ============================================================================
\section{Schritt 2: Punktoperationen}

Filter, bei denen jeder neue Pixel nur vom \textit{gleichen} Pixel im
Originalbild abh\"angt, heissen \term{Punktoperationen}.

\subsection*{Filter 2: Graustufen}

In einem Graustufenbild sind alle drei Farbkan\"ale gleich. Der einfachste
Weg: den Durchschnitt der drei Werte nehmen.

\begin{definitionbox}{Graustufen}
F\"ur jeden Pixel $(R, G, B)$:
\[ \text{grau} = \frac{R + G + B}{3}, \qquad
   \text{neuer Pixel} = (\text{grau}, \text{grau}, \text{grau}) \]
\end{definitionbox}

\begin{lstlisting}[style=pseudo]
Funktion graustufen(bild):
    Erstelle ein leeres neues_bild.
    Fuer jede Zeile des Bildes:
        Erstelle eine leere neue_zeile.
        Fuer jeden Pixel (r, g, b) der Zeile:
            Setze grau = (r + g + b) ganzzahlig durch 3.
            Haenge (grau, grau, grau) an neue_zeile an.
        Haenge neue_zeile an neues_bild an.
    Gib neues_bild zurueck.
\end{lstlisting}

\stempel{4}{Graustufen}

\subsection*{Filter 3: Helligkeit}

Um ein Bild heller zu machen, addieren wir auf jeden Farbwert eine Zahl
$h$. Achtung: die Werte m\"ussen zwischen 0 und 255 bleiben! Daf\"ur
schreiben wir eine kleine Hilfsfunktion \texttt{clamp}.

\begin{lstlisting}[style=pseudo]
Funktion clamp(wert):
    Wenn wert kleiner als 0: gib 0 zurueck.
    Wenn wert groesser als 255: gib 255 zurueck.
    Sonst: gib wert zurueck.

Funktion helligkeit(bild, h):
    Erstelle ein leeres neues_bild.
    Fuer jede Zeile des Bildes:
        Erstelle eine leere neue_zeile.
        Fuer jeden Pixel (r, g, b) der Zeile:
            Haenge (clamp(r+h), clamp(g+h), clamp(b+h)) an neue_zeile an.
        Haenge neue_zeile an neues_bild an.
    Gib neues_bild zurueck.
\end{lstlisting}

\stempel{5}{Heller, dunkler}

\begin{aufgabe}{Helligkeit von Hand}
Was kommt heraus, wenn du auf einen Pixel $(200, 100, 50)$ die Helligkeit
\begin{enumerate}
  \item $h = 50$
  \item $h = -150$
  \item $h = 100$ (ohne Clamping)
\end{enumerate}
anwendest?
\end{aufgabe}

% ============================================================================
\section{Schritt 3: Schwellwert}

Beim Schwellwert wird jeder Pixel \textit{entweder} ganz schwarz \textit{oder}
ganz weiss, je nachdem, ob seine Helligkeit \"uber oder unter einem
Schwellwert $t$ liegt.

\begin{definitionbox}{Schwellwert}
F\"ur jeden Pixel berechne zuerst den Grauwert $g$. Dann:
\[ \text{neuer Pixel} = \begin{cases}
   (255, 255, 255) & \text{falls } g \geq t \\
   (0, 0, 0) & \text{sonst}
\end{cases} \]
\end{definitionbox}

\begin{lstlisting}[style=pseudo]
Funktion schwellwert(bild, t):
    Erstelle ein leeres neues_bild.
    Fuer jede Zeile des Bildes:
        Erstelle eine leere neue_zeile.
        Fuer jeden Pixel (r, g, b) der Zeile:
            Setze grau = (r + g + b) ganzzahlig durch 3.
            Wenn grau groesser oder gleich t:
                Haenge (255, 255, 255) an neue_zeile an.
            Sonst:
                Haenge (0, 0, 0) an neue_zeile an.
        Haenge neue_zeile an neues_bild an.
    Gib neues_bild zurueck.
\end{lstlisting}

\stempel{6}{Schwarz-Weiss}

% ============================================================================
\section{Schritt 4: Nachbarpixel und Box-Blur}

Die bisherigen Filter haben jeden Pixel \textit{einzeln} bearbeitet.
Spannender wird es, wenn wir auch die \term{Nachbarpixel} mitnehmen.

\begin{definitionbox}{Box-Blur ($3\times 3$)}
F\"ur jeden Pixel: schau dir den Pixel selbst und seine 8 direkten Nachbarn
an (das $3\times 3$-Fenster). Berechne pro Farbkanal den \textbf{Durchschnitt}
aller 9 Werte. Das ist dein neuer Pixel.

\begin{center}
\begin{tikzpicture}[
  cell/.style={minimum width=10mm, minimum height=10mm, draw=maincolor!50,
    fill=maincolorvlight, font=\scriptsize},
  center/.style={cell, fill=maincolorlight, font=\scriptsize\bfseries},
]
  \node[cell] at (0, 1) {NW};
  \node[cell] at (1.2, 1) {N};
  \node[cell] at (2.4, 1) {NO};
  \node[cell] at (0, 0) {W};
  \node[center] at (1.2, 0) {Pixel};
  \node[cell] at (2.4, 0) {O};
  \node[cell] at (0, -1) {SW};
  \node[cell] at (1.2, -1) {S};
  \node[cell] at (2.4, -1) {SO};
\end{tikzpicture}
\end{center}
\end{definitionbox}

\begin{lstlisting}[style=pseudo]
Funktion box_blur(bild):
    Erstelle ein leeres neues_bild.
    Fuer jede Zeile von 1 bis hoehe-2:
        Erstelle eine leere neue_zeile.
        Fuer jede Spalte von 1 bis breite-2:
            Setze summe_r, summe_g, summe_b auf 0.
            Fuer dz in {-1, 0, 1}:
                Fuer ds in {-1, 0, 1}:
                    Lies Pixel (r, g, b) an (zeile+dz, spalte+ds).
                    summe_r = summe_r + r
                    summe_g = summe_g + g
                    summe_b = summe_b + b
            Haenge (summe_r//9, summe_g//9, summe_b//9) an neue_zeile an.
        Haenge neue_zeile an neues_bild an.
    Gib neues_bild zurueck.
\end{lstlisting}

Beachte: wir beginnen bei Zeile/Spalte~1 und h\"oren eine Zeile/Spalte vor
dem Ende auf. So haben wir immer ein vollst\"andiges $3\times 3$-Fenster.
Die \"aussersten Pixel fallen dabei weg. Das ist eine bewusste Entscheidung
und f\"ur unsere Zwecke v\"ollig in Ordnung.

\stempel{7}{Nachbarn mischen}

\begin{aufgabe}{Mehrfach blurren}
Was passiert, wenn du den Box-Blur \textit{mehrfach hintereinander} auf
dasselbe Bild anwendest? Probiere 1$\times$, 5$\times$ und 20$\times$ aus.
\end{aufgabe}

% ============================================================================
\section{Schritt 5: Bild umordnen, Spiegeln}

Bisher haben wir Pixel \textit{ver\"andert}. Jetzt ordnen wir sie um. Beim
horizontalen Spiegeln tauschen wir die linke und die rechte Seite.

\begin{lstlisting}[style=pseudo]
Funktion spiegeln(bild):
    Erstelle ein leeres neues_bild.
    Fuer jede Zeile des Bildes:
        Erstelle eine leere neue_zeile.
        Fuer jede Spalte von 0 bis breite-1:
            Haenge bild[zeile][breite-1-spalte] an neue_zeile an.
        Haenge neue_zeile an neues_bild an.
    Gib neues_bild zurueck.
\end{lstlisting}

\stempel{8}{Spiegelei}

% ============================================================================
\section{Schritt 6: Dein eigener Filter}

Jetzt bist du dran! Erfinde \textit{einen} eigenen Filter. Hier ein paar
Ideen (du darfst auch etwas ganz Eigenes machen):

\begin{itemize}
  \item \textbf{Vertikal spiegeln:} obere und untere H\"alfte tauschen.
  \item \textbf{Posterisierung:} runde jeden Kanal auf Vielfache von 64
    (Werte werden dann nur noch 0, 64, 128, 192).
  \item \textbf{Verpixelung:} ersetze jeden $k \times k$-Block durch
    die Durchschnittsfarbe dieses Blocks.
  \item \textbf{Einfache Kantenerkennung:} f\"ur jeden Pixel den Unterschied
    zum rechten Nachbarn berechnen.
  \item \textbf{Nur ein Kanal:} alle Farben ausser Rot, Gr\"un oder Blau
    auf 0 setzen.
\end{itemize}

\stempel{9}{Eigener Filter}

% ============================================================================
\section{Schritt 7: Abgabe}

Bereite die Abgabe vor:

\begin{enumerate}
  \item Pr\"ufe, dass alle Pflichtfilter laufen.
  \item Suche drei Originalbilder aus und wende je einen Filter darauf an.
    Das ergibt drei \textit{Vorher/Nachher}-Paare.
  \item Schreibe einen Bericht von 2 bis 3 A4-Seiten. Was
    war einfach? Was war schwierig? Welcher Filter hat dich \"uberrascht?
  \item Gib Code, Bilder und Bericht ab.
\end{enumerate}

\stempel{10}{Auf dem Gipfel}

\closechapter
